% /root/monodromy, 19 November 2026. Complete root and reduction/descent audits passed.
\subsection{A wreath and grid argument in degree two on every mixed product}
\label{subsec:mixed-degree-two-wreath-grid-candidate}
\status{Proved and independently checked.}
The root and the two companion reviewers each audited all five complete
proofs, their primary-source dependencies and the exact graph certificate,
including independent reruns. The result covers every fixed ordinary modulus.

Throughout this subsection $k$ is algebraically closed of characteristic
two, $X=E\times E_0$ with $E$ ordinary and $E_0$ supersingular, and
$L=\mathcal O_E(3O)\boxtimes\mathcal O_{E_0}(3O)$.
The degree-two relative unit system has rank $34$. Its entire invariant
covector space has dimension two; write $K$ for their common kernel,
of rank $32$. These ranks and covectors are the previously proved
ppav rank and nonordinary covector theorems.

\begin{lemma}[A labelled-transposition group]
\label{lem:labelled-transposition-wreath-candidate}
Let $G$ be a finite elementary abelian two-group and $N\ge3$.
Suppose $H\subset G^N\rtimes S_N$ acts transitively on
$G\times\{1,\ldots,N\}$ and is generated by labelled transpositions:
each generator interchanges two fibres, with a translation $c$ in
one direction and $-c$ in the other, and fixes all other fibres
pointwise. After changing the origins of the $N$ fibres,
\[
                   H=\ker(\Sigma:G^N\to G)\rtimes S_N.
\]
\end{lemma}
\begin{proof}
The underlying transpositions form a connected graph, by transitivity.
Choose a spanning tree from them. Successively changing the origins
along the tree makes all tree labels zero. These tree transpositions
generate the ordinary $S_N$ subgroup. If another generator exchanges
$i,j$ with label $c$, multiplying it by the unlabelled transposition
gives the translation with $c,-c$ at $i,j$. Conjugating by $S_N$
gives the same difference at any pair of indices.
Let $G_0\subset G$ be generated by all the remaining labels. The
group generated is therefore exactly
$\ker(\Sigma:G_0^N\to G_0)\rtimes S_N$.
Its orbits on $G\times\{1,\ldots,N\}$ are indexed by $G/G_0$.
Transitivity forces $G_0=G$.
\end{proof}

\begin{lemma}[Lifted intersection-point monodromy]
\label{lem:mixed-lifted-intersection-wreath-candidate}
Fix $\alpha\in\Pic^0(X)$ and a smooth sharp curve
$C\in|L\otimes\alpha|$, where sharp means two-rank nine.
Let $\widetilde C\to C$ be a connected finite etale torsor
with group $G$, a finite elementary abelian two-group.
As $D$ varies through a dense open of $|L\otimes\alpha^{-1}|$
with transverse intersection with $C$, the lifted cover of its
$N=18$ intersection points has image
\[
                     \ker(\Sigma:G^{18}\to G)\rtimes S_{18},
\]
up to choices of origins in the point fibres. The same assertion
holds after requiring $D$ to be smooth and sharp.
\end{lemma}
\begin{proof}
Use the image of $H^0(X,L\otimes\alpha^{-1})$ in
$H^0(C,(L\otimes\alpha^{-1})|_C)$ as a linear system on $C$.
It is very ample, because the line bundle on $X$ is very ample.
Its projective dimension is eight, or seven when the restriction
has its one-dimensional kernel. The original parameter space maps
onto it with integral affine fibres away from the possible centre;
this pullback does not change a finite geometric monodromy image.

We spell out the generic tangency assertion in characteristic two.
Each projection of $C$ to an elliptic factor has degree three,
so it is separable and its ramification set is finite. A degree-three
line bundle on an elliptic curve separates second jets except at
the finite set of points $P$ with that line bundle equal to
$\mathcal O(3P)$. Consequently the product line bundle separates
second jets on the surface at a general point of $C$, and hence
on $C$ itself.

For a fixed surface point the joint first-jet rank at it and a
second point is six away from the two factor fibres and the
finitely many crossed tangent incidences of
Lemma~\ref{lem:product-small-degree-jets}. The same proof
applies to the present two cubics and to their translates.
Its polar-finiteness input holds uniformly: the three partials
of $Y^2Z+XYZ+X^3+aZ^3$ are linearly independent, as seen
from the respective coefficients of $X^2,Y^2,XZ$; the
partials of $Y^2Z+YZ^2+X^3$ are $X^2,Z^2,Y^2$.
A nonzero polar quadratic therefore cannot vanish on either
entire cubic, so only finitely many tangents pass through
any prescribed point. The exceptional intersection with $C$
is finite, since both projections of $C$ are finite. Thus on
$C\times C$ the four curve first jets have rank four away from
a subset of dimension at most one. Rank two means that the two
embedded curve tangent lines coincide. At a point whose two
projection derivatives are nonzero, its tangent line contains
no other point of the Segre surface: a nonzero point on that
line has tensor rank two. The remaining points of $C$ are
finite, and each of their tangent lines meets $C$ in finitely
many points. Hence the rank-two pairs form a finite set.

The incidence of two distinct tangencies consequently has
codimension at least two in the hyperplane parameter space:
its three strata have point dimensions at most $2,1,0$ and
impose at least $4,3,2$ conditions. Higher contact at one
point also has codimension at least two, by generic second-jet
separation; at the finitely many exceptional points, the two
first-jet conditions alone give the same bound.
The tangent incidence is an irreducible projective
bundle over $C$, and its generic fibre over its image is finite.
Its image is therefore an irreducible divisor $\Delta$, whose
geometric generic member has precisely one contact of order
two and otherwise simple intersections.

The resulting unlabelled $S_{18}$ conclusion is a known
instance of the quasireflexive hyperplane-monodromy theorem
\cite[Theorem~2, p.~3; definition on p.~6; Proposition~3.1
and the proof of Theorem~2, pp.~7--8]{Entin2021}:
the jet argument above verifies its geometric hypothesis.
The labelled sum-zero wreath extension and its later gluing
application are the additional arguments proved here.

The local degree-two root cluster can be written
$z^2+bz+c=0$: after henselization the simple roots split off,
and $1,z$ form a basis of its remaining length-two algebra.
The independent constant and linear jet parameters make
$db,dc$ independent at a simple tangency, so $b,c$ are part
of a local parameter system in the coefficient base.
In particular $\Delta$ is locally $b=0$, and $c$ is not a
square at its separably closed generic residue field. The
equivalent equation $y^2+y=c/b^2$ proves that inertia interchanges
the two roots. This handles the possible purely inseparable
generic tangent-point field: no separability of the marked
tangency incidence is assumed. Since $\widetilde C\to C$ is
etale, the lifted inertia is a labelled transposition, fixing
all other point fibres.

The incidence of a point of $\widetilde C$ and a hyperplane
through its image is an irreducible projective bundle over
$\widetilde C$. Its map to the hyperplane space is finite:
no member of this embedding contains all of $C$.
It follows that the lifted-point image is transitive.
The global deck transformations commute with monodromy, so
the image is a subgroup of $G^{18}\rtimes S_{18}$.
The Galois closure has no branch divisor except $\Delta$.
Its quotient by the normal subgroup generated by divisorial
inertia is unramified in codimension one and therefore etale
over projective space by purity \cite{Stacks0BMB}; it is trivial
because geometric projective space is simply connected, as proved
in Lemma~\ref{lem:char2-defect-one-inertia}.
The image is thus generated
by these labelled transpositions. Apply the preceding lemma.

Finally smooth sharp $D$ form a nonempty dense open by the
rank theorem and the degree-one result. Restricting a finite
etale cover of a normal integral open to a dense open preserves
its connected components and monodromy image.
\end{proof}

\begin{lemma}[The nodal boundary contains all gluing translations]
\label{lem:mixed-degree-two-boundary-group-candidate}
Fix $\alpha\in\Pic^0(X)$, and vary smooth sharp transverse
$C\in|L\otimes\alpha|$, $D\in|L\otimes\alpha^{-1}|$.
Their nodal unions belong to the rank-$34$ relative unit open
inside $|L^2|$. On this boundary the relative representation
has an exact sequence
\[
 \begin{gathered}
 0\longrightarrow T\longrightarrow V\longrightarrow R\longrightarrow0,
 \qquad T=\mathbf F_2^{18}/\mathbf F_2\mathbf1,\\
 R=(H^1_{\et}(C,\mathbf F_2)\oplus
             H^1_{\et}(D,\mathbf F_2))/H^1_{\et}(X,\mathbf F_2).
 \end{gathered}
\]
Put $T^0=\{(x_i):\sum x_i=0\}/\mathbf F_2\mathbf1$,
of dimension sixteen. The boundary image contains all
translations $\operatorname{Hom}(R,T^0)$ and an actual
$\GL_8(\mathbf F_2)\times\GL_8(\mathbf F_2)$ acting naturally
on two eight-dimensional direct summands and fixing a
complement. Moreover every nonzero invariant subspace of
this boundary representation contains $T^0$.
\end{lemma}
\begin{proof}
Each normalized curve has genus ten and two-rank nine.
Their eighteen nodes add seventeen Frobenius-unit graph
directions. Thus the union has two-rank $9+9+17=35$;
the ambient rank is one, giving relative rank $34$.
The normalization sequence gives the displayed exact sequence,
with $\dim R=17$.
Write $A$ for the one-dimensional ambient unit line.
Both relative curve images are $\GL_8$ by the preceding
uniform degree-one theorem. This applies to every Picard
twist: the polarization of $L$ is an isogeny, so over $k$
every $L\otimes\alpha$ is a translate of $L$.
Since $C,D$ vary independently in a dense open of a product,
their two relative images are independent.

We give the needed group-extraction details. The image on
$H^1_{\et}(D,\mathbf F_2)$ fixes $A$ pointwise and maps
onto $\GL_8$; its unipotent kernel is either zero or the
entire natural-dual module of dimension eight. The previously
proved $H^1(\GL_8,\mathrm{nat}^*)=0$ splits the possible
complement. In either case this image is perfect and its
only nonabelian simple composition factor is $\GL_8$.
The analogous statement holds for $C$.
On $R$ the image is similarly an extension of
$L_0=\GL_8\times\GL_8$ by a submodule of
$(\mathrm{nat}_C^*\oplus\mathrm{nat}_D^*)\otimes A$.
These two simple modules are nonisomorphic; the same
first-cohomology vanishing makes the image a semidirect
product and therefore perfect.

The point-permutation image on $T$ is $S_{18}$ by the
preceding lemma. It and the image on $R$ have no common
nontrivial quotient. Indeed the latter is perfect with
nonabelian simple factors $\GL_8$, whereas the former has
only $A_{18}$ and cyclic factors. The two simple groups
are different (for example $11$ divides $|A_{18}|$ but
does not divide $|\GL_8(\mathbf F_2)|$).
Goursat's lemma gives their product as joint image.
Take the inverse image of $1\times L_0$ after choosing
the semidirect complement on $R$. Its matrix translation
kernel is an $L_0$-submodule of $\operatorname{Hom}(R,T)$.
Under $L_0$ this is a direct sum of copies of the two
natural-dual modules and the trivial module. Every quotient
has zero $H^1$: this uses the proved natural-dual vanishing
and perfectness of $L_0$ for the trivial summands.
The usual graph-of-a-cocycle argument therefore supplies
an actual block-diagonal $L_0$ in the boundary image,
fixing $T$ and a complementary ambient line in $R$.

Let $J$ be the kernel of the boundary action on $T\oplus R$.
Then $J\subset\operatorname{Hom}(R,T)$ is additive and is
an $L_0$-submodule. Fix $C$ and choose its connected maximal
elementary two-cover, whose group is
$G=H^1_{\et}(C,\mathbf F_2)^*$ of dimension nine.
The preceding lemma gives the lifted-node group
$G^{17}\rtimes S_{18}$. Its only nonabelian simple factor
is $A_{18}$. Its joint image with the moving $D$ cohomology
is their product, by the same Goursat argument: a common
quotient of the perfect $D$ image would have to have a
nonabelian simple factor, and none is shared.

Choose a pure lifted-point translation $(g_i)$ with
$\sum g_i=0$ and trivial moving-$D$ cohomology action.
It gives an element of $J$. On the injected
$H^1_{\et}(C,\mathbf F_2)\subset R$ its translation is
\[
                       \xi\longmapsto(\xi(g_i))_i
                                     \bmod\mathbf F_2\mathbf1.
\]
For precision, a class $\xi$ is a fixed etale two-torsor
on $C$. Extend it to the union by the trivial torsor on
$D$ and a chosen gluing at every node. Continuing those
gluing choices around the loop changes them by exactly
the displayed node evaluations. The two constant changes
of trivialization remove their diagonal class. This proves
the formula directly, rather than identifying an unspecified
extension class with a point-cover action.

As $(g_i)$ varies, these maps are all of
$\operatorname{Hom}(H^1_{\et}(C,\mathbf F_2),T^0)$.
The $L_0$-module $\operatorname{Hom}(R,T)$ is semisimple,
with three isotypic parts corresponding to the two
relative eight-dimensional spaces and to $A$.
An $L_0$-submodule is the direct sum of its intersections
with these parts. The preceding surjectivity therefore
puts all maps from the $C$ relative space and from $A$
to $T^0$ in $J$, with the other columns zero.
Interchanging $C,D$ supplies the remaining relative
columns. Hence $\operatorname{Hom}(R,T^0)\subset J$.

If an invariant subspace has nonzero image in $R$, these
translations force it to contain $T^0$. If it lies in
$T$, take a nonconstant representative of a nonzero
vector and apply a transposition exchanging unequal
coordinates. Their difference is the class of $e_i+e_j$.
Its $S_{18}$-orbit spans $T^0$. This also proves the last
assertion for subspaces of $T$.
\end{proof}

\begin{lemma}[Comparison at a common grid through the finite Lang part]
\label{lem:mixed-degree-two-common-grid-candidate}
Let $S$ be a nonzero sub-local system of the degree-two
relative system, and suppose $S\subset K$. Then $S=K$.
In particular the geometric representation on $K$ is irreducible.
\end{lemma}
\begin{proof}
Choose six distinct points $P_0,\ldots,P_5$ on $E$ and
six distinct points $Q_0,\ldots,Q_5$ on $E_0$, each six-point
divisor linearly equivalent to $6O$. Put
\[
 Z=\bigcup_i(\{P_i\}\times E_0)\ \cup\
   \bigcup_j(E\times\{Q_j\})\in|L^2|.
\]
Its dual graph is the complete bipartite graph $K_{6,6}$.
The graph has first Betti number $36-12+1=25$; the six
ordinary components contribute six further unit directions,
and the ambient unit line removes one. Thus the relative
stable rank at $Z$ is $30$, below the generic rank $34$.
We use only its finite Lang part, not a specialization
isomorphism of full fibres.

For any row and column three-element subsets $I,J$, group
the corresponding three vertical and three horizontal
components into $C_{I,J}$, and the complementary components
into $D_{I,J}$. Their line bundles are
$L\otimes\alpha_{I,J}$ and $L\otimes\alpha_{I,J}^{-1}$
for a suitable Picard class. This opposite twist is essential:
we do not require distinct partitions of six points to
give two triples each linearly equivalent to $3O$.
The two curves meet at the eighteen cross-edges, all simple,
and are smooth along those intersection points.
Their complete parameter systems give a boundary family
through $Z$; smooth sharp transverse pairs form a dense
open in each such family.

Here is the precise compatibility of the rank-dropping
comparison. Take the strict henselization $R_Z$ of the
local coefficient ring of $|L^2|$ at $Z$. It is a normal
local domain. The full coherent Frobenius bundle is
regular there, and its affine Lang scheme is separated,
quasi-finite and etale. Its finite part is a constant
finite etale group with special fibre the thirty-dimensional
relative unit space $W_Z$. This is the already proved
henselian finite-part construction. Its restriction to
the nonempty generic rank-$34$ open injects into $V$.

Extend $V$ and $S$ to the full coherent maximal-rank open,
including nodal unions. This does not change their images:
the original smooth open is dense in a normal integral
parameter space. For each constant section of the finite
part, membership in $S$ on that open is an open-and-closed
condition. The open is integral, so membership is either
everywhere or nowhere. These sections define one fixed
subspace $S_Z\subset W_Z$.

Pull back to the strict local parameter scheme of any
boundary family through $Z$. Its eighteen cross intersection
points remain finite etale, so its graph space $T$ and
heart $T^0$ extend as constant unit systems even at $Z$.
Their special injections into $W_Z$ agree with the
restrictions of the ambient finite Lang sections: two
sections of a separated etale scheme over a henselian
local scheme with the same closed value coincide.
The preceding boundary lemma puts $T^0$ inside $S$ at
the generic boundary point. It follows that each of these
special cross-hearts lies in the same $S_Z$.

We now compare two partitions. First take
$I_1=J_1=\{0,1,2\}$, and then
$I_2=J_2=\{0,1,3\}$.
Represent graph cohomology by edge cochains modulo vertex
coboundaries, and denote the edge $(i,j)$ by $e_{ij}$.
The second cross-heart contains
\[
                       e_{02}+e_{12},\qquad e_{34}+e_{35}.
\]
Both have even sum on its cross-edges. The first is supported
on the internal graph of $C_{I_1,J_1}$, and is nonzero
there: its evaluation on the four-cycle with rows $0,2$
and columns $1,2$ is one. The second is supported on the
internal graph of $D_{I_1,J_1}$, and its evaluation on
the four-cycle with rows $3,4$ and columns $3,4$ is one.

The two corresponding constant finite-Lang sections belong
to $S$ on the first boundary family. Projection from union
cohomology to the two component relative cohomologies is
regular and Frobenius-equivariant even at the grid. Indeed
$H^1(X,(L\otimes\alpha)^{-1})=0$ by Kunneth, so ambient
$H^1(\mathcal O_X)$ injects for every member. Connectedness,
constant arithmetic genus and cohomology-and-base-change give
the relevant coherent vector bundles. Restriction commutes
with Frobenius, as do the quotients by ambient cohomology.
The displayed graph classes survive that quotient: on each
$3\times3$ grid, ambient coherent cohomology maps injectively
to the sum of the six elliptic-component cohomologies, while
graph classes map to zero there. Hence the projections are nonzero at
$Z$ and consequently nonzero generically in that boundary
family.

The image of $S$ on the two relative component systems
therefore has nonzero projection to each of the independent
natural $\GL_8$ representations. It contains their direct
sum, of dimension sixteen. Its kernel already contains
the first cross-heart, of dimension sixteen. Thus
$\dim S\ge32=\dim K$, proving $S=K$.
\end{proof}

\begin{theorem}[Full degree-two image for every mixed product]
\label{thm:mixed-every-fixed-modulus-degree-two-candidate}
For every perfect field of characteristic two, every ordinary
elliptic curve $E$ and every supersingular elliptic curve $E_0$,
the relative system for $|L^2|$ on $E\times E_0$ has
\[
 G_{\mathrm{geom}}=G_{\mathrm{arith}}
   =\operatorname{Hom}(\mathbf F_2^2,\mathbf F_2^{32})
                         \rtimes\GL_{32}(\mathbf F_2),
\]
the full pointwise stabilizer of its two invariant covectors.
Its common covector kernel is irreducible and is contained
in every nonzero invariant subspace.
\end{theorem}
\begin{proof}
Work geometrically first. The boundary lemma shows that every
nonzero invariant subspace of $V$ contains a cross-heart
$T^0\subset K$. The common-grid lemma makes $K$ irreducible,
so every such subspace contains $K$. In particular $V$ has
no invariant vector.

In a boundary frame, the two actual natural eight-dimensional
blocks and $T^0$ lie in $K$: an invariant covector vanishes
on each of these nontrivial simple boundary constituents.
Their dimensions sum to $8+8+16=32$, so they are all of $K$.
The boundary translations include every rank-one map from
either natural block to $T^0$, acting trivially on the other
summands.

Let $N$ be the normal subgroup of the geometric image on $K$
generated by its transvections. Since that image is irreducible,
the restriction to $N$ is semisimple: the nonzero $N$-socle
is invariant under the whole group. Each isotypic multiplicity
is one, because a rank-one transvection acts nontrivially on
only one simple constituent with multiplicity one; an
isotypic component fixed by all such generators would be
$N$-fixed and hence globally invariant. Thus the simple
$N$-summands form a direct sum. The first $\GL_8$ natural
block lies in a single summand. Each rank-one map from it
to an arbitrary vector of $T^0$ forces that vector into
the same summand. The analogous maps from the second
natural block put that block in the same summand as well.
Consequently this summand is all of $K$, so $N$ is irreducible.

The transvection classification
\cite[Theorem~3.6, printed p.~22]{KantorSeress2001}, in the
already used $\mathbf F_2$ form, now gives $N=\GL(K)$ or
a group preserving a nondegenerate alternating form.
The latter alternative is impossible. The actual natural
$\GL_8$ block has no invariant bilinear form, and its
pairing with a complement fixed by that block is zero
because the natural representation has no invariant
linear functional. Such a block would lie in the radical.
Thus the image on $K$ is $\GL_{32}(\mathbf F_2)$.

The quotient $V/K$ is the constant space $\mathbf F_2^2$.
The remaining translation kernel is $K\otimes U$ for a
subspace $U\subset(\mathbf F_2^2)^*$. Modulo that kernel,
the subgroup is a graph of a cocycle. The proved vanishing
$H^1(\GL_{32},K)=0$ makes it a linear complement. If $U$
were proper, a nonzero quotient direction annihilated by
$U$ would have a fixed lift, contradicting absence of
invariant vectors. Therefore $U=(\mathbf F_2^2)^*$ and
the full displayed stabilizer occurs. The two entire
covectors descend arithmetically over every perfect base
field by the nonordinary covector theorem. Arithmetic
equality and descent of the geometric statement follow.
\end{proof}

\begin{remark}[Scope and retained checks]
\label{rem:mixed-degree-two-wreath-grid-scope-candidate}
This argument uses opposite Picard twists at the common grid
and only its finite Lang part of rank thirty. It does not
identify that special fibre with the generic rank-$34$ space.
The new uniform result concerns degree two; the fixed-modulus
all-higher-degree problem remains separate. The supporting graph
certificate is
\path{agents/monodromy/code/mixed_degree_two_grid_checks.py},
with its retained output. It checks the actual graph quotients,
cross-hearts and two nonzero projections; it does not replace
the finite-Lang comparison or geometric monodromy proof.
The author execution took $0.000195655$ seconds. Independent
root, reduction/descent and obstructions/crystals reruns took
$0.000138148$, $0.000105329$ and $0.000102624$ seconds,
respectively, preserving the original output. All geometric and
group steps have passed the complete
independent proof audits described above.
\end{remark}
